\documentclass[12pt]{article}
\usepackage{fullpage}
\usepackage[T1]{fontenc}
\usepackage[utf8]{inputenc}
\usepackage{lmodern}
\usepackage{microtype}
\usepackage{amsmath,amssymb}
\usepackage{amsthm}
\usepackage{amsfonts}
\usepackage{hyperref}

\newtheorem{theorem}{Theorem}
\newtheorem{lemma}{Lemma}
\newtheorem{proposition}{Proposition}
\theoremstyle{definition}
\newtheorem{definition}{Definition}
\newtheorem{remark}{Remark}

\DeclareMathOperator{\mult}{mult}

\title{Partition hook-length and parts generating-function identities}
\author{}
\date{}

\begin{document}
\maketitle

\section{Statement and conventions}

Throughout, a \emph{partition} of an integer $n\ge1$ is a finite non-increasing
sequence $\lambda=(\lambda_1\ge\lambda_2\ge\cdots\ge\lambda_\ell\ge1)$ of positive
integers with $\sum_i\lambda_i=n$; we write $\lambda\vdash n$ and $|\lambda|=n$. The
empty partition is the unique partition of $0$. We denote by $p(n)$ the number of
partitions of $n$, with $p(0)=1$ and $p(m)=0$ for $m<0$, and by
\begin{equation}\label{eq:Pdef}
P(x)\;=\;\prod_{k\ge1}\frac{1}{1-x^k}\;=\;\sum_{n\ge0}p(n)\,x^n
\end{equation}
Euler's partition generating function, a well-defined element of $\mathbb{Z}[[x]]$
(the product converges in the $x$-adic topology because the factor $(1-x^k)^{-1}$
equals $1+x^k+x^{2k}+\cdots\equiv 1\pmod{x^k}$, so only finitely many factors affect
each coefficient).

To a cell $u=(i,j)$ of the Young diagram of $\lambda$ (with $1\le i\le\ell$,
$1\le j\le\lambda_i$) one associates its \emph{hook length}
\begin{equation}\label{eq:hook}
h_u\;=\;(\lambda_i-j)\;+\;(\lambda'_j-i)\;+\;1,
\end{equation}
where $\lambda'_j=\#\{i:\lambda_i\ge j\}$ is the length of the $j$-th column
(conjugate partition). Each $h_u$ is a positive integer.

\paragraph{The coefficient ring.}
Fix $\alpha\in\mathbb C$ and a formal symbol $q$. As prescribed by the problem, we
work over the polynomial ring
\begin{equation}\label{eq:ring}
\mathcal R\;=\;\mathbb Z\big[Q_1,Q_2,Q_3,\dots\big],
\qquad Q_k:=q^{\,k^{\alpha}},
\end{equation}
in which the symbols $Q_k$ ($k\ge1$) are treated as \emph{algebraically independent
indeterminates}; no relation among them is assumed. All identities below are
identities of formal power series in $x$ with coefficients in $\mathcal R$, i.e.\ in
$\mathcal R[[x]]$. The two assertions to be decided are
\begin{align}
\text{(i)}\quad
&\sum_{n\ge1}x^n\sum_{\lambda\vdash n}\sum_{u\in\lambda} Q_{h_u}
\;\overset{?}{=}\;
P(x)\sum_{k\ge1} Q_k\,\frac{x^k}{1-x^k},
\label{eq:i}\\[1mm]
\text{(ii)}\quad
&\sum_{n\ge1}x^n\sum_{\lambda\vdash n}\sum_{i\ge1} Q_{\lambda_i}
\;\overset{?}{=}\;
P(x)\sum_{k\ge1} Q_k\,\frac{x^k}{1-x^k},
\label{eq:ii}
\end{align}
where in \eqref{eq:i}, $\sum_{u\in\lambda}Q_{h_u}=\sum_{u\in\lambda}q^{h_u^\alpha}$
and in \eqref{eq:ii}, $\sum_{i\ge1}Q_{\lambda_i}=\sum_{i=1}^{\ell}q^{\lambda_i^\alpha}$.

\paragraph{Main result.}
\begin{theorem}\label{thm:main}
Over the ring $\mathcal R$ of \eqref{eq:ring}:
\begin{enumerate}
\item[(ii)] Identity \eqref{eq:ii} holds in $\mathcal R[[x]]$.
\item[(i)] Identity \eqref{eq:i} is false: the coefficients of $x^2$ on the two
sides of \eqref{eq:i} differ (already in their $Q_2$-component).
\end{enumerate}
\end{theorem}

The proof occupies the rest of the paper. We first record three facts: that both
sides of \eqref{eq:i}--\eqref{eq:ii} are well-defined elements of $\mathcal R[[x]]$
(Lemma~\ref{lem:welldef}), that equality in $\mathcal R[[x]]$ is decided
coefficient-by-coefficient in the $Q_k$ (Lemma~\ref{lem:indep}), and a reduction of
the common right-hand side (Lemma~\ref{lem:rhs}).

\section{Well-definedness, independence, and the right-hand side}

\begin{lemma}[Well-definedness]\label{lem:welldef}
Each of the four series in \eqref{eq:i}--\eqref{eq:ii} is a well-defined element of
$\mathcal R[[x]]$; that is, for every $n\ge0$ the coefficient of $x^n$ is a (finite)
element of $\mathcal R$.
\end{lemma}

\begin{proof}
\emph{Left-hand sides.} For fixed $n$ there are finitely many ($p(n)$) partitions
$\lambda\vdash n$, each with finitely many cells and finitely many parts, and every
hook length $h_u$ and part $\lambda_i$ lies in $\{1,\dots,n\}$. Hence the coefficient
of $x^n$ in the left side of \eqref{eq:i} is the finite sum
$\sum_{\lambda\vdash n}\sum_{u\in\lambda}Q_{h_u}\in\mathcal R$, a polynomial in
$Q_1,\dots,Q_n$; likewise for \eqref{eq:ii}.

\emph{Right-hand side.} For each fixed $k\ge1$,
$\dfrac{x^k}{1-x^k}=\sum_{m\ge1}x^{km}\in\mathbb Z[[x]]$ starts at degree $k$, so the
sum $\sum_{k\ge1}Q_k\,\frac{x^k}{1-x^k}$ is $x$-adically summable: the coefficient of
$x^n$ receives contributions only from those $k\le n$, namely
$\sum_{k=1}^{n}Q_k\cdot[x^n]\frac{x^k}{1-x^k}\in\mathcal R$. Multiplying the
resulting element of $\mathcal R[[x]]$ by $P(x)\in\mathbb Z[[x]]$ is the usual Cauchy
product, again with finitely many terms in each degree. Hence the right-hand side
lies in $\mathcal R[[x]]$.
\end{proof}

\begin{lemma}[Coefficient extraction in the $Q_k$]\label{lem:indep}
For each $n\ge0$, the coefficient of $x^n$ in any of the series above can be written
uniquely as a finite $\mathbb Z$-linear combination
\(\sum_{k\ge1} c_k\,Q_k\) of the indeterminates $Q_k$ \emph{(}with all but finitely
many $c_k=0$\emph{)}. Consequently two such series are equal in $\mathcal R[[x]]$ if
and only if, for every $n\ge0$ and every $k\ge1$, the coefficient of $Q_k$ in the
$x^n$-coefficient agrees on both sides.
\end{lemma}

\begin{proof}
By Lemma~\ref{lem:welldef} every $x^n$-coefficient is a polynomial in the $Q_k$.
Inspecting the four series, each $x^n$-coefficient is in fact \emph{linear} in the
$Q_k$ with integer coefficients and no constant term: the left sides are sums of
single symbols $Q_{h_u}$ resp.\ $Q_{\lambda_i}$, and the right side is the Cauchy
product of $P(x)\in\mathbb Z[[x]]$ with $\sum_k Q_k\frac{x^k}{1-x^k}$, which is
$\mathbb Z[[x]]$-linear in the $Q_k$. Thus every $x^n$-coefficient lies in the free
$\mathbb Z$-module $\bigoplus_{k\ge1}\mathbb Z\,Q_k$.

Because $\mathcal R=\mathbb Z[Q_1,Q_2,\dots]$ is a polynomial ring, the monomials
$1,Q_1,Q_2,\dots$ are $\mathbb Z$-linearly independent; in particular the family
$(Q_k)_{k\ge1}$ is a $\mathbb Z$-basis of $\bigoplus_k\mathbb Z\,Q_k$. Hence the
representation $\sum_k c_kQ_k$ is unique, and equality of two elements of
$\bigoplus_k\mathbb Z\,Q_k$ is equivalent to equality of all coefficients $c_k$.
Equality in $\mathcal R[[x]]$ means equality of every $x^n$-coefficient, so the
stated coefficient-by-coefficient criterion follows. This is precisely the role of
the independence convention \eqref{eq:ring}: it converts the (a priori subtle)
identity in $q$ into a finite linear-algebra check, ruling out cancellations through
accidental relations among the actual numbers $q^{k^\alpha}$.
\end{proof}

\begin{lemma}[Right-hand side coefficients]\label{lem:rhs}
For all $n\ge1$ and $k\ge1$, the coefficient of $Q_k\,x^n$ in the common right-hand
side $\displaystyle P(x)\sum_{k\ge1}Q_k\,\frac{x^k}{1-x^k}$ equals
\begin{equation}\label{eq:Rk}
R_k(n)\;:=\;\sum_{m\ge1}p(n-km)\;=\;\sum_{m=1}^{\lfloor n/k\rfloor}p(n-km).
\end{equation}
\end{lemma}

\begin{proof}
Fix $k$. Since $\frac{x^k}{1-x^k}=\sum_{m\ge1}x^{km}$ and $P(x)=\sum_{j\ge0}p(j)x^j$,
the Cauchy product gives the contribution of the symbol $Q_k$ to the right side as
\[
Q_k\,P(x)\sum_{m\ge1}x^{km}
=Q_k\sum_{j\ge0}\sum_{m\ge1}p(j)\,x^{\,j+km}.
\]
The coefficient of $x^n$ here is $\sum_{m\ge1}p(n-km)$, where only the terms with
$km\le n$ are nonzero (as $p(j)=0$ for $j<0$), i.e.\ $1\le m\le\lfloor n/k\rfloor$.
Distinct symbols $Q_k$ contribute to distinct basis elements of
$\bigoplus_k\mathbb Z Q_k$, so by Lemma~\ref{lem:indep} the coefficient of
$Q_k\,x^n$ on the right side is exactly $R_k(n)$.
\end{proof}

\section{Proof of (ii): the parts identity}

For a partition $\lambda$ and $k\ge1$ let $\mult_k(\lambda)=\#\{i:\lambda_i=k\}$ be
the multiplicity of $k$ as a part, and set
\begin{equation}\label{eq:Nk}
N_k(n)\;:=\;\sum_{\lambda\vdash n}\mult_k(\lambda).
\end{equation}

\begin{lemma}[Linearization of the left side of \eqref{eq:ii}]\label{lem:linII}
For every $n\ge1$ the coefficient of $x^n$ in the left side of \eqref{eq:ii} equals
$\sum_{k=1}^{n} N_k(n)\,Q_k$; equivalently, the coefficient of $Q_k\,x^n$ is
$N_k(n)$.
\end{lemma}

\begin{proof}
The $x^n$-coefficient of the left side of \eqref{eq:ii} is
$\sum_{\lambda\vdash n}\sum_{i\ge1}Q_{\lambda_i}$, a finite sum since each $\lambda$
has finitely many parts. Grouping the parts of $\lambda$ by their value, for each
fixed $\lambda$ we have $\sum_{i\ge1}Q_{\lambda_i}=\sum_{k\ge1}\mult_k(\lambda)\,Q_k$
(every part equal to $k$ contributes one copy of $Q_k$). Summing over
$\lambda\vdash n$ and interchanging the two finite sums,
\[
\sum_{\lambda\vdash n}\sum_{i\ge1}Q_{\lambda_i}
=\sum_{k\ge1}\Big(\sum_{\lambda\vdash n}\mult_k(\lambda)\Big)Q_k
=\sum_{k\ge1}N_k(n)\,Q_k.
\]
Only $k\le n$ occur (a part of $\lambda\vdash n$ is at most $n$), so the sum is
finite, and by Lemma~\ref{lem:indep} the coefficient of $Q_k\,x^n$ is $N_k(n)$.
\end{proof}

\begin{lemma}[Counting parts equal to $k$]\label{lem:partcount}
For all $n\ge1$ and $k\ge1$, $\;N_k(n)=\sum_{m\ge1}p(n-km)=R_k(n)$.
\end{lemma}

\begin{proof}
Write the multiplicity as a sum of indicator functions, using
$\mult_k(\lambda)=\sum_{m\ge1}\mathbf 1[\mult_k(\lambda)\ge m]$ (a finite sum, since
$\mult_k(\lambda)\le n/k$). Then
\[
N_k(n)=\sum_{\lambda\vdash n}\sum_{m\ge1}\mathbf 1[\mult_k(\lambda)\ge m]
=\sum_{m\ge1}\#\{\lambda\vdash n:\mult_k(\lambda)\ge m\},
\]
the interchange being legitimate as all sums are finite for fixed $n$.

\emph{Bijection.} Fix $m\ge1$. Define a map $\Phi_m$ from
$\mathcal A_m:=\{\lambda\vdash n:\mult_k(\lambda)\ge m\}$ to the set of all
partitions of $n-km$ by deleting exactly $m$ of the parts equal to $k$:
\[
\Phi_m(\lambda)\;=\;\text{the partition obtained from }\lambda\text{ by removing }m
\text{ copies of the part }k.
\]
This is well-defined: if $\lambda\in\mathcal A_m$ then $\lambda$ has at least $m$
parts equal to $k$, so removing $m$ of them yields a non-increasing sequence of
positive integers summing to $n-km$, i.e.\ a partition of $n-km$ (possibly empty,
which is the partition of $0$ when $n=km$). Its inverse $\Psi_m$ adjoins $m$ parts
equal to $k$ to an arbitrary partition $\mu\vdash n-km$, inserted in the sorted order;
the result has multiplicity of $k$ at least $m$ and size $n$, hence lies in
$\mathcal A_m$. Clearly $\Phi_m\circ\Psi_m=\mathrm{id}$ and $\Psi_m\circ\Phi_m=\mathrm{id}$
(removing then re-adding, or adding then removing, exactly $m$ copies of the part
$k$ returns the original object). Therefore $\Phi_m$ is a bijection
$\mathcal A_m\to\{\mu:\mu\vdash n-km\}$, whence
\[
\#\{\lambda\vdash n:\mult_k(\lambda)\ge m\}=\#\{\mu\vdash n-km\}=p(n-km).
\]
Summing over $m\ge1$ gives $N_k(n)=\sum_{m\ge1}p(n-km)=R_k(n)$ by \eqref{eq:Rk}.
\end{proof}

\begin{proof}[Proof of Theorem~\ref{thm:main}(ii)]
By Lemma~\ref{lem:linII} the coefficient of $Q_k\,x^n$ on the left of \eqref{eq:ii}
is $N_k(n)$, and by Lemma~\ref{lem:rhs} the coefficient of $Q_k\,x^n$ on the right is
$R_k(n)$. Lemma~\ref{lem:partcount} gives $N_k(n)=R_k(n)$ for all $n,k\ge1$. By the
coefficient-extraction criterion of Lemma~\ref{lem:indep}, the two series are equal
in $\mathcal R[[x]]$. (For $n=0$ both sides vanish, since the outer sums in
\eqref{eq:ii} start at $n=1$ and $R_k(0)=0$.) This proves \eqref{eq:ii}.
\end{proof}

\section{Disproof of (i): the hook-length version fails}

For a partition $\lambda$ and $v\ge1$ let $H_v(\lambda)=\#\{u\in\lambda:h_u=v\}$ be
the number of cells of hook length $v$, and set
\begin{equation}\label{eq:Mv}
M_v(n)\;:=\;\sum_{\lambda\vdash n}H_v(\lambda).
\end{equation}

\begin{lemma}[Linearization of the left side of \eqref{eq:i}]\label{lem:linI}
For every $n\ge1$ the coefficient of $Q_v\,x^n$ in the left side of \eqref{eq:i}
equals $M_v(n)$.
\end{lemma}

\begin{proof}
The $x^n$-coefficient of the left side of \eqref{eq:i} is
$\sum_{\lambda\vdash n}\sum_{u\in\lambda}Q_{h_u}$. For fixed $\lambda$, grouping
cells by hook length gives $\sum_{u\in\lambda}Q_{h_u}=\sum_{v\ge1}H_v(\lambda)\,Q_v$;
summing over $\lambda\vdash n$ and interchanging finite sums yields
$\sum_{v\ge1}M_v(n)\,Q_v$. By Lemma~\ref{lem:indep} the coefficient of $Q_v\,x^n$ is
$M_v(n)$.
\end{proof}

\begin{proposition}[Explicit failure at $n=2$]\label{prop:counter}
$M_2(2)=2$, whereas the coefficient of $Q_2\,x^2$ on the right of \eqref{eq:i} is
$R_2(2)=1$. Hence the coefficients of $x^2$ on the two sides of \eqref{eq:i} are
distinct elements of $\mathcal R$, and \eqref{eq:i} is false.
\end{proposition}

\begin{proof}
There are exactly $p(2)=2$ partitions of $2$:
\[
\lambda=(2)\quad\text{and}\quad\lambda=(1,1).
\]
We compute hook lengths from \eqref{eq:hook}.

For $\lambda=(2)$ (one row, two columns), the conjugate is $\lambda'=(1,1)$, so
$\lambda'_1=1$, $\lambda'_2=1$. The cells are $(1,1)$ and $(1,2)$:
\[
h_{(1,1)}=(\lambda_1-1)+(\lambda'_1-1)+1=(2-1)+(1-1)+1=2,\qquad
h_{(1,2)}=(2-2)+(1-1)+1=1.
\]
Thus the hook multiset of $(2)$ is $\{2,1\}$, giving $H_2((2))=1$, $H_1((2))=1$.

For $\lambda=(1,1)$ (two rows, one column), the conjugate is $\lambda'=(2)$, so
$\lambda'_1=2$. The cells are $(1,1)$ and $(2,1)$:
\[
h_{(1,1)}=(\lambda_1-1)+(\lambda'_1-1)+1=(1-1)+(2-1)+1=2,\qquad
h_{(2,1)}=(1-1)+(2-2)+1=1.
\]
Thus the hook multiset of $(1,1)$ is $\{2,1\}$, giving $H_2((1,1))=1$, $H_1((1,1))=1$.

Summing \eqref{eq:Mv} over the two partitions,
\[
M_2(2)=H_2((2))+H_2((1,1))=1+1=2.
\]
On the other hand, by \eqref{eq:Rk},
\[
R_2(2)=\sum_{m\ge1}p(2-2m)=p(0)=1.
\]
By Lemma~\ref{lem:linI} the coefficient of $Q_2\,x^2$ on the left of \eqref{eq:i} is
$M_2(2)=2$; by Lemma~\ref{lem:rhs} the coefficient of $Q_2\,x^2$ on the right is
$R_2(2)=1$. Since the indeterminates $Q_k$ are $\mathbb Z$-linearly independent
(Lemma~\ref{lem:indep}), the $x^2$-coefficients
\[
\text{LHS:}\quad M_1(2)\,Q_1+M_2(2)\,Q_2=2Q_1+2Q_2,
\qquad
\text{RHS:}\quad R_1(2)\,Q_1+R_2(2)\,Q_2=2Q_1+Q_2
\]
are distinct elements of $\mathcal R$ (they differ in the $Q_2$-coefficient,
$2\ne1$). Here $M_1(2)=H_1((2))+H_1((1,1))=2$ and
$R_1(2)=p(1)+p(0)=1+1=2$, so the $Q_1$-parts happen to agree, but the $Q_2$-parts do
not. Therefore the two sides of \eqref{eq:i} are unequal already at degree $2$, and
\eqref{eq:i} is false.
\end{proof}

\begin{remark}[The discrepancy is genuine, not an artifact of $\alpha$]\label{rem:alpha}
The integers $M_v(n)$ and $R_v(n)$ in Lemmas~\ref{lem:rhs} and \ref{lem:linI} are
defined purely combinatorially and do not depend on $\alpha$. The failure in
Proposition~\ref{prop:counter} is the inequality $M_2(2)=2\ne1=R_2(2)$ of integers.
Under the prescribed convention \eqref{eq:ring}, where $Q_1=q^{1^\alpha}$ and
$Q_2=q^{2^\alpha}$ are independent symbols, this immediately yields
$2Q_1+2Q_2\ne2Q_1+Q_2$ for every $\alpha\in\mathbb C$. Even if one instead regarded
$q$ as a genuine variable and $\alpha$ as a fixed complex number, the two $x^2$
coefficients $2q^{1^\alpha}+2q^{2^\alpha}$ and $2q^{1^\alpha}+q^{2^\alpha}$ would
coincide only if $q^{2^\alpha}=0$, which is impossible for a nonzero formal/analytic
$q$ with $2^\alpha\ne$ a pole; hence the identity fails for every $\alpha$. The
structural reason: for $v=2$ at $n=2$, both partitions $(2)$ and $(1,1)$ contain a
cell of hook length $2$, so the \emph{total} hook count $M_2(2)=2$ exceeds the
parts count $N_2(2)=R_2(2)=1$ (only $(2)$ has a part equal to $2$).
\end{remark}

\begin{proof}[Proof of Theorem~\ref{thm:main}]
Part (ii) is established in the proof at the end of Section~3. Part (i) is
Proposition~\ref{prop:counter}, which exhibits an explicit degree-$2$ discrepancy.
\end{proof}

\section{Conclusion}

The parts identity \eqref{eq:ii} is true and follows from the classical count
$N_k(n)=\sum_{m\ge1}p(n-km)$ of parts equal to $k$ (Lemma~\ref{lem:partcount}),
matched termwise against the right-hand side $R_k(n)$ (Lemma~\ref{lem:rhs}). The
hook-length version \eqref{eq:i} is \emph{false}: the analogous statistic
$M_v(n)=\sum_{\lambda\vdash n}H_v(\lambda)$ for hook lengths does not satisfy
$M_v(n)=R_v(n)$, the smallest failure being $M_2(2)=2\ne1=R_2(2)$
(Proposition~\ref{prop:counter}).
\end{document}
